\emph{Origin and numbering.} The setup and clauses (ii) and (iii) are the setup and clauses (i) and (ii) of Theorem synthesis.positive-null-energy-gives-r-premises, which proves them for the state that Premise (N$_{+}$) supplies; here they are stated and proved for the null vacuum itself, with no premise, and the converse in (ii) (that the null vacuum is a ground state) is added. Clause (i) collects the weight and dilation statements of that result's setup, its proof of (ii)(c) and its (iv). The arguments below are moved from that proof.

Conventions. $a(f) = \int \overline{f(u)}\,\psi(u)\,du$ (antilinear in $f$), $a^{*}(f) = a(f)^{*}$; a gauge-invariant quasi-free state has symbol $A$, $\omega(a^{*}(f)a(g)) = \ip{g}{Af}$, so that its kernel is $A(u, v) = \omega(\psi^{*}(v)\psi(u))$; a flow $\tau$ is written $\tau_{t}(a(f)) = a(e^{I_{4}th}f)$ (the unit of the one-particle space is $I_{4}$; numbers are read in $\mathrm{span}\{1, I_{4}\}$), the convention of Theorem \claimref{hawking.i-state-continuum}, in which the I-state at $\beta$ has symbol $(1 + e^{\beta h})^{-1}$. Since $\nu_{a}(a(f)) = \int\overline{f(u)}\,\psi(u + a)\,du = a(f(\cdot - a))$, the one-particle generator of $\nu$ is $P = I_{4}\,d/du$ ($e^{I_{4}aP}f = f(\cdot - a)$). Since $\beta_{t}(a(g)) = \int\overline{g(u)}\,e^{\lambda/2}\psi(e^{\lambda}u)\,du = a\bigl(e^{-\lambda/2}g(e^{-\lambda}\,\cdot)\bigr)$ with $\lambda = \kappa t$, on $\varphi_{w}(u) = u^{-1/2 + I_{4}w}$ ($u > 0$) one has $e^{-\lambda/2}\varphi_{w}(e^{-\lambda}u) = e^{-I_{4}w\lambda}\varphi_{w}(u)$, so the one-particle boost generator $h$ acts on $\varphi_{w}$ as $-\kappa w$; the half-density factor $e^{\lambda/2}$ is what makes $\beta_{t}$ preserve $\{\psi(u), \psi^{*}(v)\} = \delta(u - v)$, since $\delta(e^{\lambda}(u - v)) = e^{-\lambda}\delta(u - v)$. The transverse factor $K$ is carried along as $\otimes\id_{K}$ throughout and is suppressed below.

\emph{(i)} Since $\delta(e^{\lambda}(u - v)) = e^{-\lambda}\delta(u - v)$, $e^{\lambda w_{0}}\psi(e^{\lambda}u)$ preserves $\{\psi(u), \psi^{*}(v)\} = \delta(u - v)$ exactly when $e^{2\lambda\operatorname{Re}w_{0}}e^{-\lambda} = 1$, that is $\operatorname{Re}w_{0} = \tfrac12$ (the check, in its strip function, finds $\tfrac12$ the only one among the weights $\tfrac{k}{4}$, $0 \le k \le 7$, at two rapidities). For the fields, $\beta_{t}\nu_{a}\beta_{-t}(\psi(u)) = \psi(u + e^{\kappa t}a)$ in the order that conjugates translations by the boost (the weight factors cancel), so $\beta_{t}\nu_{a}\beta_{-t} = \nu_{e^{\kappa t}a}$; the check composes the affine maps at nine rational pairs.

\emph{(ii)} Let $\omega$ be gauge-invariant and quasi-free with symbol $A$. Invariance under $\nu$ gives $e^{-I_{4}aP}Ae^{I_{4}aP} = A$ for all $a$, so $A$ is a function of $P$. A state is a ground state for $\nu$ (positive generator in the GNS representation, invariant cyclic vector) exactly when $-I_{4}\,\omega(x^{*}\delta(x)) \ge 0$ for $x$ in the domain of the derivation $\delta$ of $\nu$ (the standard characterisation of ground states, Bratteli and Robinson, Vol.~2, \S 5.3.19, cited from memory). With $x = a(g)$, $\delta(a(g)) = a(I_{4}Pg) = -I_{4}\,a(Pg)$, the condition reads $-\ip{Pg}{Ag} \ge 0$, so $PA \le 0$; with $x = a^{*}(f)$, $\delta(a^{*}(f)) = I_{4}\,a^{*}(Pf)$ and $\omega(a(f)a^{*}(g)) = \ip{f}{(1 - A)g}$, the condition reads $\ip{f}{(1 - A)Pf} \ge 0$, so $(1 - A)P \ge 0$. In the spectral representation of $P$ (in which $A$ is multiplication by $A(k) \in [0, 1]$, $0 \le A \le 1$ for any quasi-free symbol), $kA(k) \le 0$ and $(1 - A(k))k \ge 0$ force $A(k) = 0$ for $k > 0$ and $A(k) = 1$ for $k < 0$; $\{k = 0\}$ is null. So $A = \chi_{(-\infty,0)}(P)$: at most one gauge-invariant quasi-free state is a ground state for $\nu$. Since $P = I_{4}\,d/du$, $\chi_{(-\infty,0)}(P)$ is the projection onto the span of $e^{I_{4}ku}$, $k > 0$, with kernel $\tfrac{1}{2\pi}\int_{0}^{\infty}e^{I_{4}k(u - v)}\,dk = \tfrac{I_{4}}{2\pi}(u - v + I_{4}0)^{-1}$. This is the filled sea of the generator, the null vacuum $\omega_{0}$; the other ordering has kernel $\delta(u - v) - W_{0}(u, v) = -\tfrac{I_{4}}{2\pi}(u - v - I_{4}0)^{-1}$. The check confirms, in $\mathrm{span}\{1, I_{4}\}$ of $\Cl(1,3)$ and with the regulator $e^{-\epsilon k}$, that $\tfrac{1}{2\pi}\int_{0}^{\infty}e^{(I_{4}x - \epsilon)k}dk = \tfrac{1}{2\pi}(\epsilon - I_{4}x)^{-1} = \tfrac{I_{4}}{2\pi}(x + I_{4}\epsilon)^{-1}$, that the two orderings add to $\tfrac{1}{2\pi}\cdot 2\epsilon/(x^{2} + \epsilon^{2})$, whose integral over $x$ is $1$ (the anticommutator $\delta$ as $\epsilon \to 0$), and, on a rational grid of symbol values, that the two inequalities above leave exactly $A(k) = 1$ for $k < 0$ and $A(k) = 0$ for $k > 0$.

\emph{(ii), the converse.} $\omega_{0}$ is a ground state for $\nu$. Its symbol commutes with $P$, so $\omega_{0}$ is $\nu$-invariant. Its symbol is a projection, so $\omega_{0}$ is the Fock state of the splitting $\mathfrak{h} = \chi_{[0,\infty)}(P)\mathfrak{h} \oplus \chi_{(-\infty,0)}(P)\mathfrak{h}$, with particles in the first summand and holes, carried by the conjugate space, in the second; its GNS representation is that Fock representation, with $\omega_{0}$ the vacuum vector (standard for quasi-free states with projection symbol, cited from memory). The one-particle group $e^{I_{4}aP}$ commutes with the splitting, so $\nu$ is implemented by its second quantisation, whose generator is the second quantisation of $P\chi_{[0,\infty)}(P)$ on the particles and of $-P\chi_{(-\infty,0)}(P)$ on the holes (the conjugation reverses the sign of the generator), both nonnegative, and which annihilates the vacuum vector. These are the two inequalities of the necessity argument, now with equality excluded only on the null set $k = 0$.

\emph{(iii), the symbol.} From here on the complex analysis of the kernels (Mellin transform, contour rotation, continuation in the strip) is carried out in $\mathrm{span}\{1, I_{4}\}$, in which the kernels take their values. The modes beyond the edge are $\CAR(L^{2}(\R_{+}, du))$; it is $\beta$-invariant because $u \mapsto e^{\lambda}u$ preserves $\R_{+}$. The restriction of $\omega_{0}$ is gauge-invariant quasi-free with symbol $A_{+} = \chi_{+}A\chi_{+}$, kernel $\tfrac{I_{4}}{2\pi}(u - v + I_{4}0)^{-1}$ on $u, v > 0$, homogeneous of degree $-1$. Such kernels are diagonal in the Mellin basis: with $s = \tfrac12 + I_{4}w$ and $v = ux$,
\[
(A_{+}\varphi_{w})(u) = \int_{0}^{\infty}\frac{I_{4}}{2\pi}\frac{v^{s-1}\,dv}{u - v + I_{4}0} = m(w)\,\varphi_{w}(u), \qquad m(w) = \frac{I_{4}}{2\pi}\int_{0}^{\infty}\frac{x^{s-1}\,dx}{1 - x + I_{4}0} ,
\]
and the Mellin transform $f \mapsto (2\pi)^{-1/2}\int_{0}^{\infty}u^{-1/2 - I_{4}w}f(u)\,du$ is unitary from $L^{2}(\R_{+}, du)$ onto $L^{2}(\R, dw)$ (Mellin-Plancherel, cited from memory), so $A_{+}$ is multiplication by $m$. The pole of the integrand is at $x = 1 + I_{4}0$, above the axis, and the integrand is analytic in the lower half-plane and $O(|x|^{-3/2})$ at infinity and $O(|x|^{-1/2})$ at $0$; rotating the ray of integration downward through the angle $\pi$, $x = re^{-I_{4}\pi}$, crosses no singularity, and
\[
m(w) = \frac{I_{4}}{2\pi}\,e^{-I_{4}\pi s}\int_{0}^{\infty}\frac{r^{s-1}\,dr}{1 + r} = \frac{I_{4}}{2}\,\frac{e^{-I_{4}\pi s}}{\sin\pi s}
\]
by Euler's integral $\int_{0}^{\infty}r^{s-1}(1 + r)^{-1}dr = \pi/\sin\pi s$, $0 < \operatorname{Re}s < 1$ (cited from memory). With $s = \tfrac12 + I_{4}w$, $e^{-I_{4}\pi s} = -I_{4}\,e^{\pi w}$ and $\sin\pi s = \cosh\pi w$, so $m(w) = e^{\pi w}/(e^{\pi w} + e^{-\pi w}) = (1 + e^{-2\pi w})^{-1}$; equivalently $m = \tfrac{I_{4}}{2}(\cot\pi s - I_{4})$, the principal value $\pi\cot\pi s$ plus the half residue. Numerically, at $w = -0.7, 0.3, 1.25$ the rotated integral and the $I_{4}\epsilon$-regularised integral ($\epsilon = 10^{-7}$) agree with $(1 + e^{-2\pi w})^{-1}$ to $2\times 10^{-15}$ and $10^{-6}$ respectively; the opposite prescription $-I_{4}0$ gives $(1 + e^{2\pi w})^{-1}$ up to sign, the reversed orientation. Since $h\varphi_{w} = -\kappa w\,\varphi_{w}$, $(1 + e^{\beta h})^{-1}$ acts on $\varphi_{w}$ as $(1 + e^{-\beta\kappa w})^{-1}$, which is $m(w)$ for every $w$ exactly when $\beta = 2\pi/\kappa$. By Theorem \claimref{hawking.i-state-continuum}(1) (with Fact \claimref{hawking.imported-quasi-free-kms-reduces-two} lifting the two-point identity to the state), the restriction is the continuum I-state, $\beta$-KMS at $2\pi/\kappa$, and by Fact \claimref{hawking.imported-ideal-fermi-gas-uniqueness} it is the unique $\beta$-KMS state on that subalgebra. It is a single KMS state of the whole subalgebra (all $w$ and all transverse labels at once): joint, not only on marginals. Being gauge-invariant quasi-free it annihilates odd monomials, so it is even. The check verifies, with $q = e^{\pi w}$ a positive rational, that $\tfrac{I_{4}}{2}(\cot\pi s - I_{4})$ with $\cot\pi s = -I_{4}(q - q^{-1})/(q + q^{-1})$ equals $(1 + q^{-2})^{-1}$, satisfies $1 - m = q^{-2}m$, and differs from the symbols at $\beta\kappa = \pi$ and $4\pi$.

\emph{(iii), the past half.} On the generators ending on the edge, $u \le 0$, the future-directed boost is the reverse of the one on $u \ge 0$. $\CAR(L^{2}(\R_{-}))$ is invariant under $\beta$, and on it the future-directed flow is $\gamma_{t} = \beta_{-t}$. Under $u \mapsto -u$ the kernel $\tfrac{I_{4}}{2\pi}(u - v + I_{4}0)^{-1}$ on $u, v < 0$ becomes $\tfrac{-I_{4}}{2\pi}(x - y - I_{4}0)^{-1}$ on $x, y > 0$, the complex conjugate of the kernel above, so on $\varphi_{w}(x) = x^{-1/2 + I_{4}w}$ its symbol is $m(-w) = (1 + e^{2\pi w})^{-1}$; in the variable $x$ the flow $\gamma_{t}$ is $x \mapsto e^{-\kappa t}x$, whose one-particle generator acts on $\varphi_{w}$ as $+\kappa w$, so the symbol is $(1 + e^{\beta h})^{-1}$ at $\beta = 2\pi/\kappa$ again. In strip form, for $u, v < 0$, $F(\lambda) = e^{-\lambda/2}\tfrac{I_{4}}{2\pi}(e^{-\lambda}u - v + I_{4}0)^{-1}$ continued to $\lambda + I_{4}\theta$ has denominator $e^{-I_{4}\theta}e^{-\lambda}u - v$, whose imaginary part $|u|e^{-\lambda}\sin\theta$ is positive on $(0, \pi)$ and negative on $(\pi, 2\pi)$, with no zero in the open strip and the weight factor $e^{-I_{4}\pi} = -1$ at $2\pi$: KMS at $+2\pi/\kappa$ for the future-directed boost; the past-directed boost on the same half gives $-2\pi/\kappa$. The check verifies the strip signs, a tabulated value, the reversed control and $m(w) + m(-w) = 1$.

\emph{(iii), the strip and what makes $2\pi$.} The same result in the form of the KMS condition: for $u, v > 0$ put $F(\lambda) = \omega(\psi^{*}(v)\,\beta_{\lambda/\kappa}(\psi(u))) = e^{\lambda/2}\tfrac{I_{4}}{2\pi}(e^{\lambda}u - v + I_{4}0)^{-1}$. As a function of $\lambda + I_{4}\theta$, $F$ is analytic wherever $e^{I_{4}\theta}e^{\lambda}u \ne v$; for $u, v > 0$ that fails only at $\theta \in 2\pi\Z$, so $F$ is analytic in $0 < \theta < 2\pi$. At $\theta \to 0^{+}$ the denominator approaches the real axis from above, the $+I_{4}0$ of $F$; at $\theta \to 2\pi^{-}$ from below, and the weight factor has become $e^{I_{4}\pi}e^{\lambda/2}$, so $F(\lambda + 2\pi I_{4}) = e^{\lambda/2}\tfrac{-I_{4}}{2\pi}(e^{\lambda}u - v - I_{4}0)^{-1} = \omega(\beta_{\lambda/\kappa}(\psi(u))\,\psi^{*}(v))$, the second ordering computed from (ii) (its kernel is $\delta(u - v) - A(u, v) = \tfrac{-I_{4}}{2\pi}(u - v - I_{4}0)^{-1}$). This is the two-point KMS identity at rapidity period $2\pi$; smearing with test functions supported in $u > 0$ makes it rigorous, and Fact \claimref{hawking.imported-quasi-free-kms-reduces-two} lifts it. Each ingredient is needed. (a) With the opposite energy sign the prescriptions swap, the boundary value at $\theta \to 0^{+}$ no longer matches, and the strip opens to $\theta < 0$: KMS at $-2\pi/\kappa$. (b) If one point lies on the other side of the edge, $u < 0 < v$, then $e^{I_{4}\pi}e^{\lambda}u = v$ at $e^{\lambda} = v/|u|$: a pole inside the strip at imaginary rapidity $\pi$, where the continued half-generator $e^{I_{4}\theta}\R_{+}$ reaches $\R_{-}$; the full-generator algebra is not KMS (the null vacuum on it is pure). (c) The weight is not a free ingredient: $e^{\lambda w_{0}}\psi(e^{\lambda}u)$ preserves $\{\psi(u), \psi^{*}(v)\} = \delta(u - v)$ exactly when $\operatorname{Re}w_{0} = \tfrac12$, since $\delta(e^{\lambda}(u - v)) = e^{-\lambda}\delta(u - v)$; for that weight the factor at $2\pi$ is $e^{I_{4}\pi} = -1$, which carries the coefficient $I_{4}$ of one ordering to the coefficient $-I_{4}$ of the other, as the KMS identity requires. (An integer-weight bosonic field with its double-pole kernel gives the same period; the half weight is what the anticommutator forces here, not a separate reason for $2\pi$.) (d) No $\theta \in (0, 2\pi)$ is a return: $e^{I_{4}\theta}e^{\lambda}u - v$ is real only for $e^{I_{4}\theta}$ real, and at $\theta = \pi$ it is $-e^{\lambda}u - v \ne e^{\lambda}u - v$. The check verifies (a) to (d) exactly at Gaussian-rational points $e^{I_{4}\theta} = (1 - t^{2} + 2I_{4}t)/(1 + t^{2})$ of the unit circle: no zero for $u, v > 0$, the sign of the imaginary part equal to that of $t$, the zero at $\theta = \pi$ for $u < 0 < v$ with $e^{\lambda} = v/|u|$, that among the weights $\tfrac{k}{4}$ only $\tfrac12$ preserves the anticommutator ($\rho^{4w_{0}}\rho^{-2} = 1$ at $e^{\lambda/2} = \rho$) and that its factor $e^{I_{4}\pi}$ maps $I_{4}$ to $-I_{4}$, and the swapped prescription failing.
